\documentclass[10pt,a4paper]{amsart}
\usepackage[margin=19mm]{geometry}
\usepackage{amsmath,amssymb,mathtools}
\usepackage[scr=rsfs,cal=euler]{mathalfa}
\usepackage{xfrac}
\usepackage[unicode,hidelinks,bookmarks=false]{hyperref}
\newcommand{\ie}{i.e.\ }
\newcommand{\cf}{cf.\ }
\newcommand{\resp}{respectively\ }
\newcommand{\Subsection}{\subsection}
\providecommand{\nobreakdash}{\nobreak\hbox{-}}
\setlength{\emergencystretch}{2em}
\multlinegap0pt
\usepackage{amssymb}
\usepackage{enumerate}
\usepackage{mathtools}
\usepackage[normalem]{ulem}
\usepackage{cancel}
\newtheorem{lem}{Lemma}
\newcommand{\Hess}{\operatorname{Hess}}
\newcommand{\R}{{\mathbb R}}
\newcommand{\cE}{{\mathcal E}}
\newcommand{\e}{\varepsilon}
\newcommand{\supp}{\operatorname{supp}}
\newcommand{\vp}{\varphi}
\title{Partial regularity for minimizing constraint maps \\ for the Alt-Phillips energy}
\author{Rada Ziganshina}
\date{Source: arXiv:2603.29810v1 (CC BY 4.0)}
\begin{document}
\maketitle

\noindent\emph{Source and licence.} Partial regularity for minimizing constraint maps \\ for the Alt-Phillips energy. Version arXiv:2603.29810v1 (CC BY 4.0), distributed by arXiv under the Creative Commons Attribution 4.0 International licence, http://creativecommons.org/licenses/by/4.0/, as recorded in the arXiv licence metadata for this exact version and shown on the abstract page https://arxiv.org/abs/2603.29810v1. Full licence notice: \texttt{licenses/arxiv-2603.29810-CC-BY-4.0.txt}.\par\smallskip

\noindent\textit{Editorial scope.} Counted unit: the article's lem lem:dist-pde (source0.tex lines 472--481). The article's complete original proof of it is delivered with it (source0.tex lines 483--530, one \texttt{proof} environment of 48 lines). No mathematics is written, repaired or re-derived: the statement and the proof are the article's own lines, in the article's own notation, and no retained line is substituted or re-flowed.  No further source-local context is needed: every cross-reference of every retained line resolves inside the two delivered spans.  Nothing was omitted from any retained span, so the package carries no omission record.  No retained line contains a citation, so the wrapper carries no bibliography and no bibliography entry.  No retained line contains a figure environment, an image-inclusion command or drawing code, so no image is delivered and the package declares no asset dependency.  Editorial formatting only, in addition: an AMS \texttt{amsart} preamble that restates the article's own declarations and macros used by the retained lines, namely the environments \texttt{lem} on separate counters, together with the standard packages those lines need.  The retained environments print in this document as Lemma 1 in order of appearance, whereas the article prints the same items with the numbering of its own full document.  The complete native source of the article as distributed in the arXiv e-print for this exact version is bundled unchanged as \texttt{sources/197541/source0.tex}.
\begin{lem}\label{lem:dist-pde}
    Let $u\in W^{1,2}(\Omega;\overline M)$ be a minimizing constraint map for $\cE_{\lambda,\gamma}$. Suppose that $u\in C^1(U;{B_{d_0}(\partial M)})$ in a subdomain $U\subset\Omega$. Then $\rho\circ u \in C^1(U)$ and  
\begin{equation}\label{eq:dist-pde}
\begin{cases}
\Delta(\rho\circ u) = \Hess\rho_u(Du,Du) + \lambda\gamma(\rho\circ u)^{\gamma-1} & \text{in }U\cap u^{-1}(M), \\
\rho\circ u = |D(\rho\circ u)| = 0 & \text{on }U\setminus u^{-1}(M),
\end{cases}
\end{equation}
holds in the weak sense.  
\end{lem}

\begin{proof}
With $u\in C^1(U;{B_{d_0}(\partial M)})$ and $\rho\in C^\infty({B_{d_0}(\partial M)})\subset C^1({B_{d_0}(\partial M)})$, we have $\rho\circ u \in C^1(U)$. 
Since $\rho\circ u$ takes local minimum on $U\cap\partial u^{-1}(M)$, it follows that $|D(\rho\circ u)| = 0$ on $U\cap\partial u^{-1}(M)$. This proves the second line of \eqref{eq:dist-pde}. 

Since we assume that $u\in C^1(U;B_{d_0}(\partial M))\subset C(U;\R^m)$, the set $U\cap u^{-1}(M)$ is open. Moreover, since $u(U)\subset B_{d_0}(\partial M)$ and $\partial M$ is smooth, we have at least $\nu\circ u\in C^1(U)$. Therefore, given any (scalar)  function $\vp \in C_c^\infty(U\cap u^{-1}(M))$, we find $\delta > 0$ such that $\delta \leq \rho\circ u\leq d_0- \delta$ on $\supp\vp$.

Let $L > 0$ be such that $|\vp|\leq L$ on $\supp\vp$. Then for any $\e > 0$ small such that $2L\e \leq \delta$, we have 
$$0 <  \frac{1}{2}\delta  \leq \delta - \e L \leq  \rho \circ (u + \e \vp \nu\circ u)    
\leq d_0 - \delta + \e L < d_0, \qquad \hbox{in } \supp\vp.$$ 
Moreover, we also have $\rho\circ (u + \e \vp\nu\circ u) = \rho\circ u \geq 0$ in $\Omega\setminus\supp\vp$.




This shows that $u + \e\vp\nu\circ u \in W^{1,2}(\Omega;\overline M)$ and that $u + \e\vp\nu\circ u\in C^1(U;B_{d_0}(\partial M))$. Note further that since $\rho^\gamma \in C^2(M\cap B_{d_0}(\partial M))$ with $\nabla\rho^\gamma = \gamma \rho^{\gamma-1}\nu$ in $\{\rho > 0\}$, $\rho\circ u > 0$ in $u^{-1}(M)$, and $\supp\vp \Subset U\cap u^{-1}(M)$, we obtain  (using Taylor expansion)
\begin{equation}
    \label{eq:EL1}
    \begin{aligned}
        \int_\Omega \rho^\gamma\circ (u + \e \vp\nu\circ u)\,dx &= \int_\Omega \rho^\gamma\circ u \,dx  + \e \int_{U\cap u^{-1}(M)} \left(\gamma (\rho^{\gamma-1}\nu)\circ u\right)\cdot \left(\varphi \nu\circ u\right)\,dx + o(\e) \\
        &= \int_\Omega \rho^\gamma\circ u \,dx  + \e \int_{U\cap u^{-1}(M)} \gamma \varphi\rho^{\gamma-1}\circ u\,dx + o(\e),
    \end{aligned}
\end{equation}
as $\e \to 0$. Also, by the chain rule, we have $D_\alpha(\rho\circ u) = \nu^i\circ u D_\alpha u^i$ and $D_\alpha (\nu^i\circ u)D_\alpha u^i = \Hess\rho_u(D_\alpha u,D_\alpha u)$ a.e.\ in $U$, which shows that 
\begin{equation}
    \label{eq:EL2}
    \begin{aligned}
        \int_\Omega |D ( u + \e \vp\nu\circ u)|^2 \,dx &= \int_\Omega |Du|^2 \,dx \\
        &\quad + 2\e \int_{U\cap u^{-1}(M)} D(\rho\circ u)\cdot D\vp + \vp \Hess\rho_u(Du,Du)\,dx + o(\e). 
    \end{aligned}
\end{equation}
Recalling that $u + \e \vp\nu\circ u \in W^{1,2}(\Omega;\overline M)$ and that $\supp \vp \Subset U\cap u^{-1}(M)\subset\Omega$, the minimality of $u$ for $\cE_{\lambda,\gamma}$ now yields, along with \eqref{eq:EL1} and \eqref{eq:EL2}, that 
\begin{equation}
    \label{eq:EL3}
    \begin{aligned}
        \cE_{\lambda,\gamma}[u] &\leq \cE_{\lambda,\gamma}[u + \e\vp\nu\circ u] \\
        & \leq \cE_{\lambda,\gamma}[u] + 2 \e \int_{U\cap u^{-1}(M)} D(\rho\circ u)\cdot D\vp + \vp \Hess\rho_u(Du,Du) + \gamma \varphi \rho^{\gamma-1}\circ u\,dx + o(\e).
    \end{aligned}
\end{equation}
Hence, subtracting $\cE_{\lambda,\gamma}[u]$ from the left hand side and the rightmost side of \eqref{eq:EL3} and dividing the resulting inequality by $\e$ and letting $\e\to 0$ yields 
\begin{equation}
    \label{eq:EL4}
    \int_{U\cap u^{-1}(M)} D(\rho\circ u)\cdot D\vp + \vp \Hess\rho_u(Du,Du) + \gamma \varphi  \rho^{\gamma-1}\circ u\,dx \geq 0. 
\end{equation}
Finally, replacing $\vp$ with $-\vp$ turns the inequality in \eqref{eq:EL4} into an equality. As $\vp\in C_c^\infty(U\cap u^{-1}(M))$ was arbitrary, we arrive at the first line of \eqref{eq:dist-pde}. 



\end{proof}

\end{document}
